\section{加密与解密算法}

\subsection{密钥设计与混沌序列生成}

为了有效抵御选择明文攻击（CPA）与已知明文攻击（KPA），本方案设计了一种基于 SHA-512 哈希与 512 位外部密钥双重驱动的混沌参数派生机制。该机制通过明文哈希与外部密钥的全域异或、细粒度子块划分、双向对称折叠以及交叉非线性调制，实现对明文和密钥的高度敏感性。

设尺寸为 $M \times N$ 的明文图像为 $P$。首先通过安全哈希算法 SHA-512 提取明文图像的特征摘要，生成 512 位的内部哈希序列 $I = \text{SHA-512}(P)$；随后与通信双方预先商定的 512 位外部用户密钥 $K$ 进行逐位异或，得到 512 位的全局混合密钥序列 $U$：
\begin{equation}
    U = I \oplus K.
\end{equation}

为了消除单一比特扰动的局部性，系统将 512 位的混合序列 $U$ 分割为主混合段与交叉调制段：
\begin{enumerate}
    \item 前 480 位作为主混合序列，等分为 40 个 12 位的子块，记为 $[u_1, u_2, \dots, u_{40}]$，其中各子块的整数取值范围为 $[0, 2^{12}-1]$；
    \item 剩余的 32 位划分为 4 个 8 位的动态调制因子，记为 $[e_1, e_2, e_3, e_4]$，每个调制因子的整数取值范围为 $[0, 2^8-1]$。
\end{enumerate}

系统采用双端向内对称收缩的滑动结构，将全部 40 个子块完全覆盖并融合成 5 个中间归一化变量 $V_i \in [0, 1)$（$i = 1, 2, \dots, 5$）。每一组变量由序列前端的 4 个连续子块与后端的 4 个连续子块分别异或后算术加和生成：
\begin{equation}
    V_i = \frac{\left( \bigoplus_{j=1}^4 u_{4(i-1)+j} \right) + \left( \bigoplus_{j=1}^4 u_{40 - 4i + j} \right)}{2 \times 4096}, \quad i \in \{1, 2, 3, 4, 5\},
    \label{eq:Vi_v3}
\end{equation}
式中 $\oplus$ 表示按位异或运算。由于每个 12 位子块经异或后的最大取值为 $4095 < 4096$，两项算术和严格小于 $2 \times 4096 = 8192$，从而保证了 $V_i \in [0, 1)$。

随后，将归一化变量 $V_i$ 与调制因子 $e_k$ 进行正弦非线性折叠。为打破周期谐振，利用调制因子两两组合产生动态相位扰动，派生出混沌系统的控制参数 $(\mu, \nu, \alpha, \beta)$ 及初始状态 $x_0$：
\begin{equation}
\left\{
\begin{aligned}
    \mu   &= \bmod\left( \left| \sin\left( 4\pi \cdot V_1 \cdot \tilde{e}_1 \cdot \tilde{e}_4 \right) \right| \times 10^6,\; 10 \right), \\
    \nu   &= \bmod\left( \left| \sin\left( 4\pi \cdot V_2 \cdot \tilde{e}_2 \cdot \tilde{e}_3 \right) \right| \times 10^6,\; 10 \right), \\
    \alpha &= \bmod\left( \left| \sin\left( 4\pi \cdot V_3 \cdot \tilde{e}_1 \cdot \tilde{e}_2 \right) \right| \times 10^6,\; 10 \right), \\
    \beta  &= \bmod\left( \left| \sin\left( 4\pi \cdot V_4 \cdot \tilde{e}_3 \cdot \tilde{e}_4 \right) \right| \times 10^6,\; 10 \right), \\
    x_0    &= \bmod\left( \left| \sin\left( 4\pi \cdot V_5 \cdot \tilde{e}_1 \cdot \tilde{e}_3 \right) \right| \times 10^6,\; 1 \right),
\end{aligned}
\right.
\label{eq:params_v3}
\end{equation}
其中 $\bmod(z, m) = z - m\lfloor z/m \rfloor$ 为实数域连续取模运算；$\tilde{e}_k = (e_k + 0.5) / 256$ 为规避代数整数周期锁死而引入的调制系数归一化映射。最终，将所得初值 $x_0 \in (0, 1)$ 与参数代入混沌动力学系统，丢弃前 $N_0$ 次暂态迭代后生成加密所需的伪随机密钥流 $X$。


\begin{algorithm}[htbp]
\caption{SHA-512 and External Key Driven Keystream Generation}
\label{alg:keystream_complete}
\begin{algorithmic}[1]
\REQUIRE Plaintext image $P$ of size $M \times N$, 512-bit external key $K$, transient warmup count $N_0$, keystream length $L = M \times N$.
\ENSURE Dynamic parameters $(\mu, \nu, \alpha, \beta)$, initial state $x_0$, keystream sequence $X$.

\STATE $I \leftarrow \text{SHA-512}(P)$; \COMMENT{Generate 512-bit hash digest}
\STATE $W \leftarrow I \oplus K$; \COMMENT{Global bitwise XOR blending}
\STATE Split the first 480 bits of $W$ into 40 blocks of 12-bit integers: $[a_1, a_2, \dots, a_{40}]$;
\STATE Split the last 32 bits of $W$ into 4 blocks of 8-bit integers: $[b_1, b_2, b_3, b_4]$;

\FOR{$k = 1$ \TO $4$}
    \STATE $\tilde{b}_k \leftarrow (b_k + 0.5) / 256$; \COMMENT{Continuous modulation normalization}
\ENDFOR

\FOR{$i = 1$ \TO $5$}
    \STATE $S_{\text{head}} \leftarrow a_{4i-3} \oplus a_{4i-2} \oplus a_{4i-1} \oplus a_{4i}$;
    \STATE $S_{\text{tail}} \leftarrow a_{40-4i+1} \oplus a_{40-4i+2} \oplus a_{40-4i+3} \oplus a_{40-4i+4}$;
    \STATE $V_i \leftarrow \frac{S_{\text{head}} + S_{\text{tail}}}{2 \times 4096}$;
\ENDFOR

\STATE $\mu \leftarrow \bmod\left( \left| \sin(4\pi \cdot V_1 \cdot \tilde{b}_1 \cdot \tilde{b}_4) \right| \times 10^6,\; 10 \right)$;
\STATE $\nu \leftarrow \bmod\left( \left| \sin(4\pi \cdot V_2 \cdot \tilde{b}_2 \cdot \tilde{b}_3) \right| \times 10^6,\; 10 \right)$;
\STATE $\alpha \leftarrow \bmod\left( \left| \sin(4\pi \cdot V_3 \cdot \tilde{b}_1 \cdot \tilde{b}_2) \right| \times 10^6,\; 10 \right)$;
\STATE $\beta \leftarrow \bmod\left( \left| \sin(4\pi \cdot V_4 \cdot \tilde{b}_3 \cdot \tilde{b}_4) \right| \times 10^6,\; 10 \right)$;
\STATE $x_0 \leftarrow \bmod\left( \left| \sin(4\pi \cdot V_5 \cdot \tilde{b}_1 \cdot \tilde{b}_3) \right| \times 10^6,\; 1 \right)$;

\STATE $x \leftarrow x_0$;
\FOR{$t = 1$ \TO $N_0 + L$}
    \STATE $x \leftarrow \left| \sin\left( (5 + 3\mu) \left( 1 - \nu \cdot x \cdot \sin(15\pi x (1 - x)) \right) \right) \right|$;
    \IF{$t > N_0$}
        \STATE $X[t - N_0] \leftarrow x$;
    \ENDIF
\ENDFOR
\RETURN $\mu, \nu, \alpha, \beta, x_0, X$;
\end{algorithmic}
\end{algorithm}